
    %%%%%%%%%%%%%%%%%%%%%%%%%%%%%%%%%%%%%%%%%%%%%%%%%%%%%%%%%%%%%%%%%%%%%%%%%%%%%%%%
    %%% Classe do documento
    \documentclass[12pt, addpoints]{exam}
    \UseRawInputEncoding

    %%%%%%%%%%%%%%%%%%%%%%%%%%%%%%%%%%%%%%%%%%%%%%%%%%%%%%%%%%%%%%%%%%%%%%%%%%%%%%%%
    % preambulo.tex
    %
    % Arquivo de configuração de packages para uso o LaTeX, conforme minhas
    % preferências e modelos pessoais.
    %
    % Para maiores informações, visite:
    %    https://github.com/abrantesasf/latex
    %
    % NÃO ALTERE SE NÃO SOUBER O QUE ESTÁ FAZENDO!
    %%%%%%%%%%%%%%%%%%%%%%%%%%%%%%%%%%%%%%%%%%%%%%%%%%%%%%%%%%%%%%%%%%%%%%%%%%%%%%%%


    %%%%%%%%%%%%%%%%%%%%%%%%%%%%%%%%%%%%%%%%%%%%%%%%%%%%%%%%%%%%%%%%%%%%%%%%%%%%%%%%
    %%% Estruturas de controle:
    \input{static/utils/controles}

    %%%%%%%%%%%%%%%%%%%%%%%%%%%%%%%%%%%%%%%%%%%%%%%%%%%%%%%%%%%%%%%%%%%%%%%%%%%%%%%%
    %%% Compilação condicional em PDF:
    \input{static/utils/pdf}

    %%%%%%%%%%%%%%%%%%%%%%%%%%%%%%%%%%%%%%%%%%%%%%%%%%%%%%%%%%%%%%%%%%%%%%%%%%%%%%%%
    %%% Configurações de layout da página:
    \input{static/utils/layout}

    %%%%%%%%%%%%%%%%%%%%%%%%%%%%%%%%%%%%%%%%%%%%%%%%%%%%%%%%%%%%%%%%%%%%%%%%%%%%%%%%
    %%% Configurações para autores:
    \input{static/utils/autores}

    %%%%%%%%%%%%%%%%%%%%%%%%%%%%%%%%%%%%%%%%%%%%%%%%%%%%%%%%%%%%%%%%%%%%%%%%%%%%%%%%
    %%% Cabeçalho e rodapé:
    %%% Atenção! É MUITO DIFÍCIL ajustar apropriadamente os running headers
    %%% e running footers da primeira vez. Você pode confiar no LaTeX e deixar que
    %%% ele faça o ajuste automaticamente ou pode quebrar a cabeça e
    %%% tentar melhorar algo que o LaTeX já faz muito bem, mesmo que não fique
    %%% exatamente do jeito que você quer. O que você prefere? Se quiser ajustar
    %%% manualmente, descomente a linha abaixo e faça as configurações no arquivo
    %%% "utils/headings.tex".
    %\input{static/utils/headings}

    %%%%%%%%%%%%%%%%%%%%%%%%%%%%%%%%%%%%%%%%%%%%%%%%%%%%%%%%%%%%%%%%%%%%%%%%%%%%%%%%
    %%% Configuração para as fontes do documento:
    %%% Se você quiser usar fontes próprias ou do sistema, precisará ajustar as
    %%% configurações no arquivo "utils/fontes.tex".
    %%% ATENÇÃO: por padrão eu utilizo XeLaTeX com fontes proprietárias projetadas
    %%% por Matthew Butterick, adquiridas em https://mbtype.com, que não podem ser
    %%% distribuídas devido à licença de utilização. Se você usar XeLaTeX, você
    %%% DEVERÁ OBRIGATORIAMENTE ajustar as fontes no arquivo "utils/fontes.tex".
    \input{static/utils/fontes}

    %%%%%%%%%%%%%%%%%%%%%%%%%%%%%%%%%%%%%%%%%%%%%%%%%%%%%%%%%%%%%%%%%%%%%%%%%%%%%%%%
    %%% Sumário:
    \input{static/utils/toc}

    %%%%%%%%%%%%%%%%%%%%%%%%%%%%%%%%%%%%%%%%%%%%%%%%%%%%%%%%%%%%%%%%%%%%%%%%%%%%%%%%
    %%% Matemática:
    \input{static/utils/matematica}

    %%%%%%%%%%%%%%%%%%%%%%%%%%%%%%%%%%%%%%%%%%%%%%%%%%%%%%%%%%%%%%%%%%%%%%%%%%%%%%%%
    %%% Notas de rodapé:
    \input{static/utils/footnotes}

    %%%%%%%%%%%%%%%%%%%%%%%%%%%%%%%%%%%%%%%%%%%%%%%%%%%%%%%%%%%%%%%%%%%%%%%%%%%%%%%%
    %%% Referências cruzadas, links, citações:
    \input{static/utils/crossrefs}

    %%%%%%%%%%%%%%%%%%%%%%%%%%%%%%%%%%%%%%%%%%%%%%%%%%%%%%%%%%%%%%%%%%%%%%%%%%%%%%%%
    %%% Configurações para os hiperlinks em PDF:
    \input{static/utils/pdfconfig}

    %%%%%%%%%%%%%%%%%%%%%%%%%%%%%%%%%%%%%%%%%%%%%%%%%%%%%%%%%%%%%%%%%%%%%%%%%%%%%%%%
    %%% Glossário e índice remissivo:
    \input{static/utils/glossindex}

    %%%%%%%%%%%%%%%%%%%%%%%%%%%%%%%%%%%%%%%%%%%%%%%%%%%%%%%%%%%%%%%%%%%%%%%%%%%%%%%%
    %%% Referências bibliográficas:
    \input{static/utils/bibliografia}

    %%%%%%%%%%%%%%%%%%%%%%%%%%%%%%%%%%%%%%%%%%%%%%%%%%%%%%%%%%%%%%%%%%%%%%%%%%%%%%%%
    %%% Cores:
    \input{static/utils/cores}

    %%%%%%%%%%%%%%%%%%%%%%%%%%%%%%%%%%%%%%%%%%%%%%%%%%%%%%%%%%%%%%%%%%%%%%%%%%%%%%%%
    %%% Computação:
    \input{static/utils/computacao}

    %%%%%%%%%%%%%%%%%%%%%%%%%%%%%%%%%%%%%%%%%%%%%%%%%%%%%%%%%%%%%%%%%%%%%%%%%%%%%%%%
    %%% Gráficos:
    \input{static/utils/graficos}

    %%%%%%%%%%%%%%%%%%%%%%%%%%%%%%%%%%%%%%%%%%%%%%%%%%%%%%%%%%%%%%%%%%%%%%%%%%%%%%%%
    %%% Ícones:
    \input{static/utils/icones}

    %%%%%%%%%%%%%%%%%%%%%%%%%%%%%%%%%%%%%%%%%%%%%%%%%%%%%%%%%%%%%%%%%%%%%%%%%%%%%%%%
    %%% Blocos:
    \input{static/utils/blocos}

    %%%%%%%%%%%%%%%%%%%%%%%%%%%%%%%%%%%%%%%%%%%%%%%%%%%%%%%%%%%%%%%%%%%%%%%%%%%%%%%%
    %%% Boxes coloridos:
    \input{static/utils/colorbox}

    %%%%%%%%%%%%%%%%%%%%%%%%%%%%%%%%%%%%%%%%%%%%%%%%%%%%%%%%%%%%%%%%%%%%%%%%%%%%%%%%
    %%% Tabelas:
    \input{static/utils/tabelas}

    %%%%%%%%%%%%%%%%%%%%%%%%%%%%%%%%%%%%%%%%%%%%%%%%%%%%%%%%%%%%%%%%%%%%%%%%%%%%%%%%
    %%% Ambientes de listas:
    \input{static/utils/listas}

    %%%%%%%%%%%%%%%%%%%%%%%%%%%%%%%%%%%%%%%%%%%%%%%%%%%%%%%%%%%%%%%%%%%%%%%%%%%%%%%%
    %%% Ambientes floats e similares:
    \input{static/utils/floats}

    %%%%%%%%%%%%%%%%%%%%%%%%%%%%%%%%%%%%%%%%%%%%%%%%%%%%%%%%%%%%%%%%%%%%%%%%%%%%%%%%
    %%% Ajustes para a classe EXAM:
    \input{static/utils/exam}

    %%%%%%%%%%%%%%%%%%%%%%%%%%%%%%%%%%%%%%%%%%%%%%%%%%%%%%%%%%%%%%%%%%%%%%%%%%%%%%%%
    %%% Ajustes para textos:
    \input{static/utils/texto}

    %%%%%%%%%%%%%%%%%%%%%%%%%%%%%%%%%%%%%%%%%%%%%%%%%%%%%%%%%%%%%%%%%%%%%%%%%%%%%%%%
    %%% Ambientes, aliases, comandos e customizações em geral para serem
    %%% utilizadas nos documentos. Defina aqui qualquer customização específica
    %%% para seus documentos!
    \input{static/utils/customizacoes}

    %%%%%%%%%%%%%%%%%%%%%%%%%%%%%%%%%%%%%%%%%%%%%%%%%%%%%%%%%%%%%%%%%%%%%%%%%%%%%%%%
    %%% Ajuste do layout, espaçamento de linhas e etc.:
    \geometry{a4paper, portrait, left=2.0cm, right=2.0cm, top=2.0cm, bottom=2.0cm}
    \singlespacing

    %%%%%%%%%%%%%%%%%%%%%%%%%%%%%%%%%%%%%%%%%%%%%%%%%%%%%%%%%%%%%%%%%%%%%%%%%%%%%%%%
    %%% Configurações de cabeçalho e rodapé (não use fancyhdr pois ocorre conflito):
    \pagestyle{headandfoot}
    \runningheadrule
    \firstpageheader{}{}{}
    \runningheader{Sem disciplina}{}{Avaliação}
    \firstpagefooter{}{}{Página \thepage\ de \numpages}
    \runningfooter{}{}{Página \thepage\ de \numpages}
    \runningfootrule

    %%%%%%%%%%%%%%%%%%%%%%%%%%%%%%%%%%%%%%%%%%%%%%%%%%%%%%%%%%%%%%%%%%%%%%%%%%%%%%%%
    %%% Configurações para as propriedades do PDF:
    \hypersetup{
    hidelinks,           % Comente para web, descomente para impressão
    colorlinks=true,      % True para web, False para impressão
    pdftitle=sss: Avaliação,
    pdfauthor=bbb,
    pdfsubject=Sem disciplina,
    pdfkeywords=['Sem tag'],
    pdfinfo={
        CreationDate={}, % Ex.: D:AAAAMMDDHH24MISS
        ModDate={}       % Ex.: D:AAAAMMDDHH24MISS
    }
    }
    \input{static/utils/pdfconfig}

    %%%%%%%%%%%%%%%%%%%%%%%%%%%%%%%%%%%%%%%%%%%%%%%%%%%%%%%%%%%%%%%%%%%%%%%%%%%%%%%%
    %%% Começa o documento
    \begin{document}

    %%%%%%%%%%%%%%%%%%%%%%%%%%%%%%%%%%%%%%%%%%%%%%%%%%%%%%%%%%%%%%%%%%%%%%%%%%%%%%%%
    %%% Folha de instruções padronizadas da UVV para a prova
    \pdfbookmark[1]{Instruções}{instrucoes}
    \begin{figure}[H]
        \begin{center}
            \includegraphics[scale=0.3]{{static/imgs/logo-uvv.png}}
        \end{center}	
    \end{figure}

    \vspace{-1cm}

    \begin{table}[H]
        \centering
        \begin{tabular}{|llll|}
            \hline
            \multicolumn{2}{|l|}{\textbf{Disciplina:} Sem disciplina} & 
            \multicolumn{1}{l|}{\multirow{3}{*}{\textbf{\makecell{Nota:\\ \,\\ \,}}}} & 
            \multirow{3}{*}{\textbf{\makecell{Coordenador:\\ \,\\ \,}}} \\ 
            \cline{1-2}
            \multicolumn{2}{|l|}{\textbf{Professor:} bbb \hspace{4.72cm} } & 
            \multicolumn{1}{l|}{} & \\ 
            \cline{1-2}
            \multicolumn{2}{|l|}{\textbf{Aluno:}} & 
            \multicolumn{1}{l|}{} & \\ 
            \hline
            \multicolumn{1}{|l|}{\textbf{Turma:} \hspace{6.3cm} } & 
            \multicolumn{1}{l|}{\textbf{Semestre:}} & 
            \multicolumn{2}{l|}{\textbf{Valor:} 10 pontos} \\ 
            \hline
            \multicolumn{1}{|l|}{\textbf{Data:}} & 
            \multicolumn{3}{l|}{\textbf{Avaliação:}} \\ 
            \hline
        \end{tabular}
    \end{table}

    \begin{center}
        \fbox{\setlength\fboxsep{0.3cm} \fbox{\parbox{15.4cm}{
            \begin{center}
                \textbf{INSTRUÇÕES PARA A REALIZAÇÃO DESTA AVALIAÇÃO:}
            \end{center}
            \begin{itemize}[leftmargin=*]
                \item Esta é a primeira avaliação da disciplina Sem disciplina, 
                    e engloba o conteúdo referente Sem tag.
                \item Esta avaliação contém 10 questões objetivas, \textbf{todas obrigatórias}.
                \item \textbf{Leia atentamente cada questão!} As questões objetivas terão uma,
                    e apenas uma única, resposta a ser marcada.
                \item \textbf{Desligue o celular} antes de começar. A avaliação é
                    \textbf{individual} e \textbf{sem consulta}.
                \item As questões podem ser respondidas, na prova, com lápis ou caneta. Ao final
                    da prova \textbf{{VOCÊ DEVERÁ PREENCHER O GABARITO COM CANETA
                    DE COR PRETA}} \textbf{\underline{OU AZUL}} para que seja feita a correção
                    automatizada através da leitura do padrão de respostas marcado no
                    gabarito.
                \item \textbf{CUIDADO ao preencher o gabarito} pois eles são \textbf{INDIVIDUAIS
                    e NOMEADOS} (cada estudante tem um gabarito próprio específico, com um
                    código de identificação único). \textbf{Questões rasuradas são anuladas}
                    pelo software de leitura do gabarito.
                \item Ao final da prova, \textbf{DEVOLVA A PROVA E O GABARITO} para o professor.
                \item Siga todas as normas de \textbf{Integridade Acadêmica} da disciplina.
                    Alunos que forem flagrados com qualquer espécie de ``cola'' ou trocando
                    informações com outros alunos terão suas avaliações recolhidas, as notas
                    zeradas, e a situação será encaminhada para a coordenação para a aplicação
                    das penalidades previstas pela universidade.
                \item Com 90 minutos de avaliação você terá tempo de até 9,min por questão,
                    em média.
                \item Boa avaliação!
            \end{itemize}
            \vspace{2cm}
        }}}
    \end{center}
    \newpage

    %%%%%%%%%%%%%%%%%%%%%%%%%%%%%%%%%%%%%%%%%%%%%%%%%%%%%%%%%%%%%%%%%%%%%%%%%%%%%%%%
    %%% Inicia o bloco de questões:
    \begin{questions}

    \newpage
    \question
Assinale a alternativa que contém a ordem correta das cláusulas SQL:

\begin{choices}
\choice Nenhuma das alternativas acima
\choice SELECT, FROM, JOIN, WHERE, GROUP BY, ORDER BY, HAVING
\choice SELECT, FROM, WHERE, JOIN, HAVING, GROUP BY, ORDER BY
\choice SELECT, FROM, JOIN, WHERE, GROUP BY, HAVING, ORDER BY
\choice SELECT, FROM, WHERE, GROUP BY, JOIN, HAVING, ORDER BY
\end{choices}

\question
Qual das seguintes opções \textbf{não é} uma vantagem do uso de um banco de
dados sobre o sistema de arquivos convencional?

\begin{choices}
\choice Estruturação de dados
\choice Controle de redundância e inconsistência
\choice Eficiência
\choice Controle de acesso e visualização de dados
\choice Independência de dados
\end{choices}

\question
Professor Mac Sperto propôs o seguinte algoritmo de ordenação, chamado de Super Merge, similar ao \textit{merge sort}: divida o vetor em 4 partes do mesmo tamanho (ao invés de 2, como é feito no \textit{merge sort}). Ordene recursivamente cada uma das partes e depois intercale-as por um procedimento semelhante ao procedimento de intercalação do \textit{merge sort}. Qual das alternativas abaixo é verdadeira?
\begin{choices}
\choice Nenhuma das afirmações acima está correta
\choice Super Merge está correto, mas consome tempo menor que \( O(\text{merge sort}) \)
\choice Super Merge está correto, mas consome tempo maior que \( O(\text{merge sort}) \)
\choice Super Merge está correto, mas consome tempo \( O(\text{merge sort}) \)
\choice Super Merge não está correto. Não é possível ordenar quebrando o vetor em 4 partes
\end{choices}

\question
\begin{figure}[H]
    \centering
    \includegraphics[width=0.5\linewidth]{static/imgs/defaul.png}
    \caption{Enter Caption}
\end{figure}
Se \( O = (0, 0, 0) \); \( A = (2, 4, 1) \); \( B = (3, 1, 1) \) e \( C = (1, 3, 5) \) então o volume do sólido acima é

\begin{choices}
\choice 35/2
\choice 21
\choice 44
\choice 30
\choice 35
\end{choices}

\question
(ENADE 2017) No modelo relacional de dados, uma junção entre tabelas cria uma
outra tabela derivada de duas ou mais tabelas de acordo com os critérios
especificados para a junção, de modo similar às regras da teoria dos
conjuntos. Nos conjuntos a seguir, as áreas hachuradas representam os resultados
das consultas SQL em que ``id' é uma chave primária e ``fk' é uma chave
estrangeira:
\begin{figure}[H]
  \centering
  \includegraphics[width=0.5\linewidth]{static/imgs/defaul.png}
\end{figure}
Assinale a opção em que a declaração SQL é equivalente ao conjunto apresentado a
seguir:
\begin{figure}[H]
  \centering
  \includegraphics[width=0.5\linewidth]{static/imgs/defaul.png}
\end{figure}

\begin{choices}
\choice \begin{verbatim}SELECT *
FROM tabela1 t1
RIGHT JOIN tabela2 t2 ON (t1.id = t2.fk)
WHERE t2.fk <> NULL;
\end{verbatim}

\choice \begin{verbatim}SELECT *
FROM tabela1 t1 WHERE EXISTS (SELECT 1
                              FROM tabela2 t2
                              WHERE t2.id = t1.fk);
\end{verbatim}

\choice \begin{verbatim}SELECT *
FROM tabela1 t1
RIGHT JOIN tabela2 t2 ON (t1.id = t2.fk)
WHERE t1.id IS NULL;
\end{verbatim}

\choice \begin{verbatim}SELECT *
FROM tabela1 t1
INNER JOIN tabela2 t2 ON (t1.id = t2.fk);
\end{verbatim}

\choice \begin{verbatim}SELECT *
FROM tabela1 t1
WHERE NOT EXISTS (SELECT 1
                  FROM tabela1 t2
                  WHERE t2.id = tk.fk);
\end{verbatim}

\end{choices}

\question
São modelos de dados legados, pré-relacionais, exceto:

\begin{choices}
\choice Hierárquico
\choice Nenhuma das alternativas anteriores
\choice Flat file
\choice Rede
\choice Objeto
\end{choices}

\question
Engenharia de Software inclui um grande número de teorias, conceitos, modelos, técnicas e métodos. Analise as seguintes definições.
\begin{enumerate}[label=\Roman*.]
    \item O processo de inferir ou reconstruir um modelo de mais alto nível (projeto ou especificação) a partir de um documento de mais baixo nível (tipicamente um código fonte);
    \item Capacidade de modificação de um software (ou de um de seus componentes) após sua entrega ao cliente visando corrigir falhas, expandir a funcionalidade, modificar a performance ou outros atributos em resposta a novos requisitos do usuário ou mesmo ser adaptado a alguma mudança do ambiente de execução (plataforma, p.ex);
    \item Modelo estabelecido pelo Software Engineering Institute (SEI) que propõe níveis de competência organizacional relacionados à qualidade do processo de desenvolvimento de software;
\end{enumerate}
Estas definições correspondem respectivamente aos seguintes termos:
\begin{choices}
\choice engenharia reversa, reparabilidade, Team Software Process (TSP)
\choice reengenharia, manutenibilidade, Capability Maturity Model (CMM)
\choice reengenharia, evolutibilidade, Personal Software Process (PSP)
\choice engenharia reversa, manutenibilidade, Capability Maturity Model (CMM)
\choice engenharia reversa, reparabilidade, Team Software Process (TSP)
\end{choices}

\question
Qual o valor do atributo \( E.val \) após a análise da expressão "4 / 2 / 2" para o esquema de tradução a seguir?

\begin{itemize}
    \item \( E \rightarrow T / E1 \ \{ E.val = T.val / E1.val \} \)
    \item \( E \rightarrow T \ \{ E.val = T.val \} \)
    \item \( T \rightarrow \text{digito} \ \{ T.val = \text{val(digito)} \} \)
\end{itemize}

\begin{choices}
\choice 8
\choice 2
\choice 4
\choice 1
\choice 3
\end{choices}

\question
Considere as seguintes tabelas em uma base de dados relacional (chaves primárias sublinhadas):\\
Departamento (\underline{CodDepto}, NomeDepto)\\
Empregado (\underline{CodEmp}, NomeEmp, CodDepto)\\
Considere as seguintes restrições de integridade sobre esta base de dados relacional:\\
- Empregado.CodDepto é sempre diferente de NULL\\
- Empregado.CodDepto é chave estrangeira da tabela Departamento com cláusulas ON DELETE RESTRICT e ON UPDATE RESTRICT\\
Qual das seguintes validações não é especificada por estas restrições de integridade:
\begin{choices}
\choice Sempre que uma linha for excluída de Departamento, deve ser garantido que o valor de Departamento.CodDepto não aparece na coluna Empregado.CodDepto.
\choice Sempre que o valor de Departamento.CodDepto for alterado, deve ser garantido que não há uma linha com o antigo valor de Departamento.CodDepto na coluna Empregado.CodDepto.
\choice Sempre que uma nova linha for inserida em Empregado, deve ser garantido que o valor de Empregado.CodDepto aparece na coluna Departamento.CodDepto.
\choice Sempre que uma nova linha for inserida em Departamento, deve ser garantido que o valor de Departamento.CodDepto aparece na coluna Empregado.CodDepto.
\choice Sempre que o valor de Empregado.CodDepto for alterado, deve ser garantido que o novo valor de Empregado.CodDepto aparece em Departamento.CodDepto.
\end{choices}

\question
Considere as seguintes afirmativas sobre as linguagens usadas para análise sintática:
\begin{enumerate}[label=\Roman*.]
    \item A classe LL(1) não aceita linguagens com produções que apresentem recursões diretas à esquerda (ex. \(L \rightarrow La\)) mas aceita linguagens com recursões indiretas (ex. \(L \rightarrow Ra, R \rightarrow Lb\)).
    \item A linguagem LR(1) reconhece a mesma classe de linguagens que LALR(1).
    \item A linguagem SLR(1) reconhece uma classe de linguagens maior que LR(0).
\end{enumerate}
Selecione a afirmativa correta:
\begin{choices}
\choice As afirmativas I e II são verdadeiras
\choice Apenas a afirmativa III é verdadeira
\choice As afirmativas I e III são verdadeiras
\choice As afirmativas II e III são verdadeiras
\choice As afirmativas I e III são falsas
\end{choices}



    %%%%%%%%%%%%%%%%%%%%%%%%%%%%%%%%%%%%%%%%%%%%%%%%%%%%%%%%%%%%%%%%%%%%%%%%%%%%%%%%
    %%% Finaliza o bloco de questões:
    \end{questions}
    \end{document}
    